\chapterhead{2}{Big Shadow, Small Shadow}{What decides how large a shadow is}\label{ch2}

Hold your hand close to a torch and the shadow on the wall is enormous. Pull your hand away and it shrinks. Your hand did not grow or shrink. The \emph{shadow} did.

Light travels in straight lines, so you can predict the size of a shadow by drawing the lines. Draw a line from the light past the top of the object. Where it hits the wall is the top of the shadow.

\begin{center}
\begin{tikzpicture}[x=0.1cm,y=0.1cm]
  % wall at x=60, torch at x=0, object height 10
  \draw[line width=1.2pt] (60,-4) -- (60,36);
  \fill[sun!80,draw=ink] (0,0) circle (3);
  \node[font=\sffamily\scriptsize] at (0,-7) {light};
  \draw[line width=1pt,dashed,color=ink!40] (0,0) -- (60,0);
  % object at 20 (height 10)
  \draw[line width=3pt,color=ink] (20,0) -- (20,10);
  \draw[line width=1pt,color=tomato] (0,0) -- (60,30);
  \draw[<->,line width=0.9pt] (63,0) -- (63,30) node[midway,right,font=\sffamily\scriptsize] {30};
  % object at 40
  \draw[line width=3pt,color=sky!80!black] (40,0) -- (40,10);
  \draw[line width=1pt,color=sky!80!black] (0,0) -- (60,15);
  \draw[<->,line width=0.9pt] (56,0) -- (56,15) node[midway,left,font=\sffamily\scriptsize] {15};
  \node[font=\sffamily\scriptsize] at (20,-7) {object at 20};
  \node[font=\sffamily\scriptsize] at (40,-7) {object at 40};
  \node[font=\sffamily\scriptsize,anchor=north] at (60,-6) {wall at 60};
\end{tikzpicture}
\end{center}

The red line passes the top of the object at 20~cm from the light. It reaches the wall at 30~cm up. The blue line passes the top of the same object at 40~cm from the light. It reaches the wall at only 15~cm. The closer the object is to the light, the bigger its shadow.

\begin{tryit}{Shadow puppet test}
Question: \textbf{does moving a puppet closer to the torch make its shadow taller?} Use the fair-test habits from Book~2.
\begin{enumerate}[leftmargin=1.6em]
\item Cut a card puppet 10~cm tall. Tape it to a stick so it stands up.
\item Put a torch on a stack of books so it sits at the puppet's height. Tape a white sheet to the wall, 60~cm from the torch.
\item Hold the puppet 20~cm from the torch, then 30, 40 and 50~cm. Keep the torch and the wall where they are.
\item At each distance, measure the shadow's height on the wall, to the nearest cm. Do each distance twice.
\end{enumerate}
\par\medskip
\renewcommand{\arraystretch}{1.6}
\begin{tabular}{|c|c|c|c|}
\hline
\textbf{Distance from light (cm)} & \textbf{Shadow, try 1 (cm)} & \textbf{Shadow, try 2 (cm)} & \textbf{Average (cm)}\\ \hline
20 & & & \\ \hline
30 & & & \\ \hline
40 & & & \\ \hline
50 & & & \\ \hline
\end{tabular}
\end{tryit}
\adult{Use a torch or phone light. Do not use a hot bulb that someone could touch.}

\subhead{What stays the same, and what changes}
\begin{center}
\sffamily\small\renewcommand{\arraystretch}{1.4}
\begin{tabular}{@{}p{3.4cm}p{7.6cm}@{}}
\toprule
\textbf{I change} & how far the puppet is from the light: 20, 30, 40 or 50~cm\\
\textbf{I measure} & the height of the shadow on the wall, nearest cm\\
\textbf{I keep the same} & the puppet (10~cm tall), the torch, the wall's distance (60~cm), the torch height\\
\bottomrule
\end{tabular}
\end{center}

\begin{multix}[hmm]\multixbot[hmm]\tcblower
I have a plan. I will move the puppet closer to the light \emph{and} move the wall further away, to get an even bigger shadow. Then I will say the puppet position made it bigger. Is that a fair test?
\end{multix}

\subhead{Drawing predicts, measuring checks}
Before you measure, you can predict by drawing. The drawing above says the shadow should be 30, 20, 15 and 12~cm at distances 20, 30, 40 and 50~cm. Here are practice numbers from a made-up run.

\begin{worked}{Shadow puppet results (practice numbers, not from a real test)}
\begin{center}\sffamily\small\renewcommand{\arraystretch}{1.3}
\begin{tabular}{@{}cccc@{}}
\textbf{Distance (cm)} & \textbf{Drawing says (cm)} & \textbf{Measured (cm)} & \textbf{Off by (cm)}\\
20 & 30 & 29 & 1\\
30 & 20 & 20 & 0\\
40 & 15 & 16 & 1\\
50 & 12 & 12 & 0
\end{tabular}
\end{center}
\end{worked}

The drawing was within 1~cm each time. That is a good sign for the drawing.

\clearpage

\begin{center}
\begin{tikzpicture}[x=0.065cm,y=0.1cm]
  \draw[line width=1pt] (0,0) -- (0,36) (0,0) -- (62,0);
  \foreach \x in {20,30,40,50}{\draw (\x,-1) -- (\x,1); \node[font=\sffamily\scriptsize,below=2pt] at (\x,-1) {\x};}
  \foreach \y in {10,20,30}{\draw (-1,\y) -- (1,\y); \node[font=\sffamily\scriptsize,left=2pt] at (-1,\y) {\y};}
  \node[font=\sffamily\scriptsize] at (36,-8) {distance from light (cm)};
  \node[font=\sffamily\scriptsize,rotate=90] at (-10,18) {shadow height (cm)};
  \draw[line width=1.2pt,color=sky,dashed] (20,30) -- (30,20) -- (40,15) -- (50,12);
  \foreach \x/\y in {20/29,30/20,40/16,50/12}{\fill[tomato,draw=ink] (\x,\y) circle (1.2);}
\end{tikzpicture}
\end{center}
The dashed line is the drawing's prediction and the dots are the measurements. They sit almost on top of each other.

\trust{We drew a prediction first, then measured the shadow heights, changing only the puppet's distance from the light and doing each distance twice.}
{When a drawing made from straight lines predicts the measured heights to within 1~cm at four distances, the straight-line idea is working. Changing only the distance means the distance, not the torch or wall, is the likely reason for the change.}
{That the drawing works for fuzzy lights, curved objects or very big shadows. It also does not show \emph{why} light goes straight. Four distances and one puppet are a small test.}

\subhead{A pattern in the numbers}
Look at the last column of the table in the worked box: the shadow gets smaller as the puppet moves away, but not by the same amount each time. From 20 to 30~cm it drops by 10~cm, from 30 to 40 it drops by 5, and from 40 to 50 it drops by 3. Equal steps in distance do not give equal steps in shadow height.

\begin{keepbox}
We predicted shadows with straight-line drawings. \textbf{We kept} \wline[3.4cm] \textbf{and we set aside} \wline[3.4cm].
\end{keepbox}

\begin{puzzle}{Predict the shadow}
Use the drawing method: the object is 10~cm tall and the wall is 60~cm from the light.\\
(a) The object is 15~cm from the light. How tall is the shadow? \wline[2cm]~cm\\
(b) The object is 60~cm from the light, right against the wall. How tall is the shadow? \wline[2cm]~cm\\
Hint: shadow height $=10\times60\div$ distance.
\end{puzzle}

\begin{puzzle}{Spot the unfair test}
Mara tests how the distance from the light changes a shadow. She moves the puppet closer to the torch, and she also moves the torch to a lower shelf. The shadow gets taller.\\
(a) What did Mara mean to change? \wline[4cm]\\
(b) What else changed? \wline[4cm]\\
(c) How would you fix the test? \wline[4cm]
\end{puzzle}

\begin{puzzle}{Check the numbers}
A student measured these shadows for a 10~cm puppet at 20, 30, 40, 50~cm: 30, 20, 15, 11~cm.\\
(a) What does the drawing predict for 50~cm? \wline[1.5cm]\\
(b) How far off was the 11~cm reading? \wline[1.5cm]\\
(c) Does this change your trust in the drawing? Why or why not?\\
\rule{\linewidth}{0.5pt}
\end{puzzle}

\begin{puzzle}{Bigger or smaller?}
Write B (bigger) or S (smaller) for the shadow:
\begin{description}[leftmargin=1.6em,itemsep=0.3em,style=sameline]
\item[(a)] You move your hand toward the torch. \wline[1cm]
\item[(b)] You move the wall closer to the torch and the hand stays still. \wline[1cm]
\item[(c)] You use a taller puppet at the same distance. \wline[1cm]
\end{description}
\end{puzzle}
\clearpage

\subhead{More shadow sums}
Use the same rule: shadow height $=$ object height $\times$ wall distance $\div$ object distance.

\begin{puzzle}{Your hand}
Your hand is 18~cm tall. The wall is 60~cm from the torch.\\
(a) The hand is 20~cm from the torch. Shadow: \wline[1.6cm]~cm\\
(b) The hand is 40~cm from the torch. Shadow: \wline[1.6cm]~cm
\end{puzzle}

\begin{puzzle}{Two puppets}
A 10~cm puppet and a 20~cm puppet are both 30~cm from the torch, wall at 60~cm.\\
Shadows: \wline[1.6cm]~cm and \wline[1.6cm]~cm. How many times taller is the second shadow? \wline[1.6cm]
\end{puzzle}

\begin{puzzle}{Move the wall}
A 10~cm puppet is 20~cm from the torch. The wall is moved from 60~cm away to 80~cm away.\\
New shadow height: \wline[1.6cm]~cm. Did it grow or shrink? \wline[2cm]
\end{puzzle}

\begin{puzzle}{Read the graph}
From the graph earlier in this chapter, about how tall was the shadow at 30~cm? \wline[1.6cm]~cm\\
At what distance was it about 15~cm? \wline[1.6cm]~cm
\end{puzzle}

\begin{puzzle}{Explain to a friend}
In one or two sentences, explain why the shadow of your hand grows when you bring it closer to the torch.\\
\rule{\linewidth}{0.5pt}\\[4pt]\rule{\linewidth}{0.5pt}
\end{puzzle}
\clearpage
